
\subsubsection{2.1 Code LLM Benchmarks}

HumanEval \citep{Chen2021} established the pass@k metric for code generation evaluation,
using 164 hand-written Python problems with unit tests. MBPP \citep{Austin2021} provided
974 Python programming tasks ranging from simple to challenging. Both benchmarks evaluate
functional correctness via execution.

\subsubsection{2.2 Heuristic-Only Corpus Filtering}

StarCoder [Li et al., 2023a] and StarCoder 2 \citep{Lozhkov2024} established the
standard pipeline for open-source code corpus construction from The Stack [Kocetkov et al.,
2022]: near-deduplication, PII redaction, and heuristic quality filtering (line length,
alphanumeric fraction, encoding, XML/HTML ratio). These filters remove obvious junk but do
not apply execution-based quality signals. StarCoder achieves 40\% HumanEval pass@1 at
15.5B parameters.

Heuristic filtering cannot distinguish syntactically valid from invalid code or functionally
correct from incorrect code.

\subsubsection{2.3 Quality Curation via Language Models}

phi-1 \citep{Gunasekar2023} demonstrated that 7B quality tokens curated by GPT-4 yield
50.6\% HumanEval pass@1 at 1.3B parameters --- establishing that quality can outweigh quantity
at equal token budget. phi-1.5 [Li et al., 2023b] extended this to reasoning tasks.
GPT-4-based quality curation, however, is not reproducible or scalable for open-source use.

\subsubsection{2.4 Execution-Based Data Selection}

EffiCoder \citep{Zeng2024} applies execution-based sample selection to instruction-tuning
data (prompt-solution pairs), achieving +13pp HumanEval pass@1 on Qwen2.5-Coder-7B.
OpenCodeInstruct [2025] provides large-scale instruction-tuning data with execution feedback.
Token Cleaning [2025] applies token-level influence filtering to SFT data.

These works operate on instruction-tuning pairs with labeled test cases, not raw corpus code.
Raw corpus execution filtering requires intrinsic signals (compile(), doctest execution), and
the feasibility of these signals at corpus scale has not been measured prior to this work.

\subsubsection{2.5 Summary}

\begin{table}[htbp]
\centering
\resizebox{\columnwidth}{!}{%
\begin{tabular}{llll}
\toprule
Prior Work & Data Type & Execution Gate & Feasibility Measured \\
\midrule
StarCoder [Li et al., 2023a] & Raw corpus & Heuristic only & No \\
phi-1 [Gunasekar et al., 2023] & GPT-4 curated & No & No \\
EffiCoder [Zeng et al., 2024] & Instruction pairs & Execution pass rate & No \\
**This work** & **Raw corpus** & **compile() + doctest** & **Yes** \\
\bottomrule
\end{tabular}
}
\end{table}

